%% Explicitly derived calculus for the declared negative semantic interface.
\section{A Derived Calculus for the Negative Free-Logic Semantics}
\label{sec:p08-negative-free-logic-calculus}

The negative row of Table~\ref{tab:p08-free-logic-variants} fixes a nonempty
outer domain, a nonempty existing inner domain, total ordinary functions,
strict predication and identity, and inner-domain quantifiers.  Those clauses
are the starting point here.  The following calculus is an explicitly
constructed translation calculus with a complete adequacy proof; it is not
claimed as the original calculus of \cite{Lambert1967,Bencivenga1986,Lehmann2002}.
The function-free sequent presentation in
\cite[\S\S1.2, 1.3, 3.4]{GratzlPavlovicNolt2025} is a separate source
presentation.  Its non-inclusivity rule and atomic restrictions help expose
why domain and term policies must be fixed before transferring a rule.

\subsection{Language, parameters, and the partial-description boundary}

\begin{hbdefinition}{def:p08-negative-language}{The ordinary language and the strict description extension}
Fix a finitary first-order signature, with constants and total ordinary
function symbols, and a countably infinite supply of variables.  The
ordinary language has terms generated by variables, constants and function
application; its formulas have predicate atoms, strict equality \(r=t\),
the existence atom \(E!(t)\), falsum, classical Boolean connectives and
\(\forall,\exists\).  Ordinary predicates may have arity zero.  Bound
variables are alpha-renamed before capture-avoiding simultaneous
substitution.

The \emph{strict description extension} additionally permits the mutually
generated finite term constructor \(\iota xA\).  Its value is the unique
existing object satisfying \(A\), if there is exactly one, and is undefined
otherwise.  An ordinary function applied to extended terms has its total
outer-domain value when all arguments are defined, and is undefined if any
argument is undefined.  This last clause is the explicit strict lifting of
the total ordinary function, not a requirement that defined function
arguments or results be existing objects.

The ordinary language supplies the base treatment below.  The
extended result uses precisely the stated strict lifting.  This is the
explicit function/description completion of the selected negative
convention, not new historical wording attributed to the listed sources.
On the fragment
where description terms occur as complete atomic arguments and no
undefined argument is passed into a function, this extension agrees directly
with the primitive-description clauses following
Definition~\ref{def:p08-negative-free-logic}.  The strict-lifting clause
completes the nested function/description syntax explicitly; it is not
silently read into an unspecified non-strict extension.
\end{hbdefinition}

\begin{hbdefinition}{def:p08-negative-models}{Negative models and existing-parameter contexts}
A negative model has nonempty sets \(D\subseteq U\), where \(U\) is the
outer domain and \(D\) the existing domain.  Every ordinary constant denotes
in \(U\), every \(n\)-ary ordinary function is a total map \(U^n\to U\),
and every ordinary predicate denotes a subset of \(D^n\).  Variable
assignments \(g\) take values in \(U\); bound quantifiers change an
assignment only to a value in \(D\).  In the description extension terms
have the partial evaluation of Definition~\ref{def:p08-negative-language}.
An atom is false if any of its complete argument terms is undefined or
nonexisting.  Otherwise predicate atoms test their predicate extension,
strict equality tests equality of the two existing values, and \(E!(t)\)
tests existence of the value.  Nullary atoms simply have their assigned
Boolean truth value.  All Boolean operations are two-valued classical
operations.

An \emph{existing-parameter context} \(X\) is a set of variable names
required to satisfy \(g(x)\in D\).  Write
\(\Gamma\models_{-,X}A\) when every such model and every outer assignment
satisfying these parameter requirements and every member of \(\Gamma\)
also satisfies \(A\).  Other free variables may take any outer value.
Thus \(X=\varnothing\) is outer-assignment consequence; taking
\(X\) to contain all free variables in \(\Gamma\cup\{A\}\) gives
inner-assignment consequence.  Sentence consequence is independent of this
choice.  The context is a condition on assignments, not another set of
object-language premises or a restriction on ordinary constants.
\end{hbdefinition}

The selected negative convention uses outer assignments for free variables,
so its ordinary unindexed consequence is the \(X=\varnothing\) case.
The inner-parameter comparison is named separately through \(X\); it does
not replace that selected convention.  Quantifiers still assign their bound
variables only to inner-domain objects.

The declared inner and outer domains are both nonempty.  The existence of an
inner object will therefore be present in the translation theory.  The
translation does not allow an empty inner domain and then describe the result
as completeness for the existing nonempty-domain semantics.

\subsection{A finite relational translation of partial term values}

\begin{hbdefinition}{def:p08-negative-translation}{Value formulas and formula translation}
The target is ordinary classical first-order logic with equality, over the
same ordinary function and predicate symbols and one fresh unary predicate
\(E\).  Its domain is the entire outer domain.  Its equality is ordinary
outer identity, not the source strict-equality atom.  Use an additional
infinite supply of target variables for fresh bound auxiliaries.

Define simultaneously the target value formula \(V_t(z)\) and formula
translation \(\delta A\).  The output variable \(z\) is fresh for \(t\),
and all auxiliary variables are fresh for the surrounding expressions.
Enumerate the auxiliary supply once and take its first available fresh
names at each finite recursive clause; this fixes the translation rather
than leaving a free choice of names in its definition.  Every displayed
substitution avoids capture.  For variables, constants and ordinary function
application, including extended argument terms when admitted,
\[
\begin{aligned}
 V_x(z)&=(z=x),& V_c(z)&=(z=c),\\
 V_{f(t_1,\ldots,t_n)}(z)
 &=\exists u_1\cdots\exists u_n
   \left(\bigwedge_{i=1}^{n}V_{t_i}(u_i)
               \land z=f(u_1,\ldots,u_n)\right).
\end{aligned}
\]
For \(n=0\), this is just \(z=f()\), with no existential variables or
conjunctive requirements.  In the description extension,
\[
 V_{\iota xA}(z)=E(z)\land(\delta A)[z/x]\land
   \forall u\left((E(u)\land(\delta A)[u/x])\to u=z\right).
\]
The binding variable of a description is renamed before this clause, so its
substitutions refer to its free argument occurrences in the body.

For an ordinary predicate of positive arity, put
\[
 \delta P(t_1,\ldots,t_n)=
 \exists u_1\cdots\exists u_n
 \left(\bigwedge_{i=1}^{n}(V_{t_i}(u_i)\land E(u_i))
                    \land P(u_1,\ldots,u_n)\right).
\]
For nullary \(P\), put \(\delta P=P\).  The distinguished atoms translate as
\[
\begin{aligned}
 \delta E!(t)&=\exists u\,(V_t(u)\land E(u)),\\
 \delta(r=t)&=\exists u\exists v\,
       (V_r(u)\land V_t(v)\land E(u)\land E(v)\land u=v).
\end{aligned}
\]
Translation fixes falsum and commutes with every Boolean connective.  Its
quantifier clauses are
\[
 \delta(\forall xA)=\forall x(E(x)\to\delta A),\qquad
 \delta(\exists xA)=\exists x(E(x)\land\delta A).
\]
In each clause the recursive calls are on proper syntax-tree constituents.
The simultaneous recursion therefore gives finite target formulas, including
for nested descriptions; there is no circular definition of a term by the
adequacy theorem.  Put
\[
 \Theta=\{\exists zE(z)\},\qquad
 \mathcal E_X=\{E(x):x\in X\},\qquad
 \delta[\Gamma]=\{\delta B:B\in\Gamma\}.
\]
Here \(\Theta\) is a closed target theory asserting inner nonemptiness,
and \(\mathcal E_X\) is the target assignment-guard set.
\end{hbdefinition}

The fresh predicate \(E\) is not an ordinary source predicate supplied with
an arbitrary extension: it encodes precisely the existing domain.  The target
ordinary predicate extensions may contain outer tuples, but the guards in
\(\delta P\) make those tuples irrelevant.  No claim that the two model
classes are literally identical is needed.

\begin{hbtheorem}{thm:p08-negative-model-bridge}{proved=annex:p8.negative.model-bridge}{Value and truth correspondence in both model directions}
Every negative model has a classical target encoding with domain \(U\),
\(E=D\), and the same ordinary functions and predicates.  Conversely every
classical target model of \(\Theta\) gives a negative model by taking
\(D=E\)'s extension and restricting each ordinary predicate to \(D^n\).
For either direction and every outer assignment \(g\):
\begin{enumerate}
\item \(V_t(z)\) is true with \(z\) assigned \(a\in U\) exactly when
      \(t\) is defined with value \(a\);
\item at most one value satisfies \(V_t\), and every ordinary term has a
      value;
\item the source formula \(A\) is true exactly when the target formula
      \(\delta A\) is true under the same free-variable assignment.
\end{enumerate}
These assertions hold for the ordinary language, and for the strict
description extension under its stated function-lifting clause.
Consequently, for every \(X,\Gamma,A\),
\[
 \Gamma\models_{-,X}A
 \quad\Longleftrightarrow\quad
 \Theta\cup\mathcal E_X\cup\delta[\Gamma]
                         \models_{\mathsf{FOL}}\delta A.
\]
\end{hbtheorem}

\begin{hblemma}{thm:p08-negative-substitution}{proved=annex:p8.negative.substitution}{Substitution when the replacing term has a value}
If \(t\) is defined under \(g\) with value \(a\), capture-avoiding
substitution satisfies
\[
 \operatorname{val}_g(r[t/x])=
       \operatorname{val}_{g[x\mapsto a]}(r),\qquad
 M,g\models A[t/x]\ \Longleftrightarrow\ M,g[x\mapsto a]\models A.
\]
Equality of partial values includes the case where both sides are undefined.
The formulas need not be atomic; the assertion includes quantifiers and
description bodies.  It requires definedness of \(t\), not existence; the
guarded quantifier rules additionally require its value to belong to \(D\).
\end{hblemma}

\subsection{Certificates, guarded rules, and strong completeness}

\begin{hbdefinition}{def:p08-negative-calculus}{The translation-certificate calculus}
Write \(\Gamma\vdash^{\mathrm{tr}}_{-,X}A\) when there is a finite proof
on the exact classical Sheet~\ref{sheet:nd-fol} of \(\delta A\), whose
open assumptions are drawn from
\(\Theta\cup\mathcal E_X\cup\delta[\Gamma]\).  Such a proof is a
\emph{certificate} of the negative judgment.  It uses the target equality,
target quantifiers and target grammar throughout; raw source formulas are
not passed unchanged to classical quantifier or reflexivity rules.
For an arbitrary premise set, each certificate still uses only finitely
many premises and parameter guards.  The certificate records that finite
premise and parameter support.  Checking its proof structure is the ordinary
finite target proof-checking problem; it does not decide membership in an
arbitrary noneffectively supplied \(\Gamma\) or \(X\).  Semantic entailment
is not an inference-rule side condition.

The new sheet's certificate row is the defining rule.  Its remaining rows
are derived native-language interfaces proved below, and are not an assertion
that this menu without certificates is an original complete natural-
deduction presentation.  At a native interface a finite set of hypotheses
records the permitted open assumptions; unused hypotheses may be present.
The classical proof inside a certificate retains the discharge convention
of the registered target sheet.
\end{hbdefinition}

\showcalculus{negative-free-logic}

\begin{hbtheorem}{thm:p08-negative-adequacy}{proved=annex:p8.negative.adequacy}{Strong adequacy of the derived negative calculus}
For every premise set \(\Gamma\), existing-parameter context \(X\), and
formula \(A\) of the stated ordinary language,
\[
 \Gamma\vdash^{\mathrm{tr}}_{-,X}A
 \quad\Longleftrightarrow\quad \Gamma\models_{-,X}A.
\]
The same equivalence holds for the declared strict description extension.
For countable signatures this uses no Axiom of Choice; for arbitrary
signatures it uses the ambient Choice assumption of
Theorem~\ref{thm:fol-completeness}.  Soundness needs neither assumption.
If a consequence holds, some finite \(\Gamma_0\subseteq\Gamma\) and
finite \(X_0\subseteq X\) already give the corresponding consequence.
Thus these exact consequence relations are compact.  For sentences, this is
adequacy for the nonempty inner/outer negative semantics independently of
the free-variable convention.  In particular \(X=\varnothing\) proves
strong completeness for the selected outer-assignment consequence,
including its strict partial-description extension.
\end{hbtheorem}

The classical dependencies are the full natural-deduction soundness result,
Theorem~\ref{thm:fol-soundness}, and strong completeness in
Theorem~\ref{thm:fol-completeness} or its countable-signature form,
Corollary~\ref{cor:fol-completeness-countable}.  The preceding model
and truth proof establishes the transfer to the original domain, atomic,
identity, and term policies.  Merely calling the relation a translation
calculus would not establish that transfer.

\begin{hbtheorem}{thm:p08-negative-native-rules}{proved=annex:p8.negative.rules}{Admissibility of every displayed native interface}
Every derived row on Sheet~\ref{sheet:negative-free-logic} preserves
\(\models_{-,X}\) and is admissible for
\(\vdash^{\mathrm{tr}}_{-,X}\).  The guarded quantifier steps require
\(E!(t)\); the generalization/elimination eigenvariable is fresh for
\(X\), the other open hypotheses, and the conclusion as specified on
the sheet.  Equality reflexivity requires existence, and equality
replacement uses strict identity, hence the same existing value.
An atom yields existence of each of its complete argument terms, not of
arbitrary proper subterms.  The nonempty-domain rule discharges the
fresh existing witness only because \(D\ne\varnothing\).
\end{hbtheorem}

The theorem is proved after adequacy, so the native-interface admissibility
proof does not serve as an assumption in the completeness proof.  Classical
Boolean reasoning is available because the entire negative semantics is
two-valued; it does not remove the guards from quantifier or equality rows.

\begin{hbcorollary}{thm:p08-negative-equality-existence}{proved=annex:p8.negative.equality-existence}{Existence, strict self-identity, and existing witnesses}
For every term \(t\), including a possibly failed description,
\[
 E!(t)\ \dashv\vdash^{\mathrm{tr}}_{-,X}\ t=t
 \ \dashv\vdash^{\mathrm{tr}}_{-,X}\ \exists x(x=t),
\]
where the witness variable is fresh for \(t\).  Also
\(\vdash^{\mathrm{tr}}_{-,X}\exists xE!(x)\).
The existence premise cannot be dropped to obtain unrestricted
\(\vdash^{\mathrm{tr}}_{-,X}t=t\).
\end{hbcorollary}

Here the displayed double turnstiles abbreviate two separate finite
derivability implications between the displayed single formulas.  They do
not equate a formula with a model property or with a target equality.

\begin{hbexample}{def:p08-negative-worked-guards}{An existing witness and an unknown constant}
From \(\forall xP(x)\) and \(E!(c)\), guarded universal elimination
derives \(P(c)\), and guarded existential introduction derives
\(\exists xP(x)\).  Independently, from a positive-arity atom \(P(c)\)
its atomic-existence interface supplies \(E!(c)\), so the same existential
conclusion is available.  By contrast, \(\neg E!(c)\) is not a source
atom: it may be true for a nonexisting constant and does not supply
\(E!(c)\).  Existentially generalizing that negative formula without a
guard is unsound.
\end{hbexample}

\begin{hblemma}{thm:p08-negative-certificate-example}{proved=annex:p8.negative.certificate-example}{A target certificate of guarded universal elimination}
For an ordinary constant \(c\) and unary ordinary predicate \(P\),
\[
 \{\forall xP(x),E!(c)\}\vdash^{\mathrm{tr}}_{-,\varnothing}P(c).
\]
The reciprocal annex gives an actual ND-FOL certificate of its translated
conclusion, including every existential value witness, rather than treating
the raw free-logic premises as classical formulas.
\end{hblemma}

\subsection{Separating cases and the limits of the transfer}

\begin{hbtheorem}{thm:p08-negative-boundaries}{proved=annex:p8.negative.boundaries}{Explicit failures of stronger or unspecified conventions}
The first four separating cases use nonempty outer and inner domains.
The last case exhibits failure after changing the inner-domain convention.
\begin{enumerate}
\item Let \(U=\{0,1\}\), \(D=\{0\}\), and \(c=1\).
      Then \(\forall xE!(x)\) and \(\neg E!(c)\) are true, while
      \(E!(c)\), \(c=c\), and \(\exists x\neg E!(x)\) are false.
      Thus unguarded universal instantiation, unguarded existential
      generalization, and unconditional reflexivity fail.
\item In the same model with a unary predicate true exactly of \(0\),
      let a total ordinary function satisfy \(f(1)=0\).  Then
      \(P(f(c))\) is true and \(E!(c)\) false.  Atomic truth does not
      imply existence of every syntactically occurring subterm.  No closure
      of \(D\) under ordinary functions, or strictness with respect to
      nonexistence, was imposed in the declared semantics.
\item With these domains, the free formula \(E!(x)\) is valid if every
      free-variable assignment is inner-valued, but fails under the outer
      assignment \(g(x)=1\).  The two conventions are distinguished by
      \(X\); neither is inferred merely from the domain clause.
\item In the description extension, \(\iota x\bot\) is undefined.
      An atom having this term as a complete argument, possibly after
      ordinary function application propagates its undefinedness, is false;
      its negation is true, and the term's strict self-identity false.
      If \(D\) has two objects, \(\iota xE!(x)\)
      is also undefined; failed uniqueness has the same negative atomic
      policy as failed existence.
\item An empty existing domain would falsify \(\exists xE!(x)\) and
      invalidate the sheet's nonempty-domain rule.  It is outside the model
      class just proved complete; inclusive negative logic is a different
      domain convention.
\end{enumerate}
\end{hbtheorem}

The strict description extension propagates undefinedness through function
application, but not nonexistence: a defined nonexisting argument may have
an existing function value.  These two distinctions are needed simultaneously.
The preceding proof covers arbitrary nested strict descriptions.  It does
not assign a value to an application with an undefined input under a
separately chosen non-strict policy.  Likewise, unindexed open-formula
consequence must specify which free parameters are inner-valued; the two
typed presentations here keep that choice visible.

The construction proves strong adequacy, finite support, and admissibility
of the guarded interfaces.  It asserts neither a cut-elimination theorem
for a native sequent calculus nor a normalization procedure for every
translation certificate.  The different original free-logic presentations,
including their own description and function policies, retain their separate
source obligations \cite{Lambert1967,Bencivenga1986,Lehmann2002}.

\index{negative free logic!derived translation calculus}
\index{free logic!existence guards}\index{strict identity}
\index{undefined term}\index{nonexisting object}\index{free variable!domain convention}
